% Generated draft: requires mathematical and source review.
\documentclass[11pt]{article}
\usepackage[margin=25mm]{geometry}
\usepackage{fontspec}
\setmainfont{Latin Modern Roman}
\usepackage{amsmath,amssymb,mathtools,unicode-math}
\setmathfont{Latin Modern Math}
\usepackage{enumitem,needspace,xurl,hyperref}
\setlength{\parindent}{0pt}
\setlength{\parskip}{5pt}
\setlength{\emergencystretch}{3em}
\newcommand{\N}{\mathbb{N}}
\newcommand{\Z}{\mathbb{Z}}
\newcommand{\Q}{\mathbb{Q}}
\newcommand{\R}{\mathbb{R}}
\newcommand{\C}{\mathbb{C}}
\newcommand{\F}{\mathbb{F}}
\newcommand{\A}{\mathbb{A}}
\newcommand{\PP}{\mathbb{P}}
\newcommand{\E}{\mathbb{E}}
\newcommand{\eps}{\varepsilon}
\newcommand{\dd}{\,\mathrm{d}}
\begin{document}
\section*{UnsolvedMath ID: 30000508}
\textbf{Source problem code:} catalogue entry
\subsection{Definitions and mathematical statement}
For the period-matrix locus $H^+(U)$, is the imaginary axis the image of a Teichm\"uller geodesic?

The supplied assessment places this question in a special infinite-genus Teichm\"uller construction. The parameter $U$, locus $H^+(U)$, underlying surface, marking conventions, and period map are not defined in the extraction. In particular, no finite-dimensional Siegel-space model or finite-genus period theorem is substituted.

If $\Phi$ denotes the intended period map from the relevant Teichm\"uller space, the requested existence would amount to a geodesic $\gamma:I\to\mathcal T$ whose set of period images is exactly the source's imaginary axis:
\[
 \Phi(\gamma(I))=\{\text{points on the imaginary axis intended in }H^+(U)\}.
\]
This formula records the source's image requirement, not a definition of the missing axis or map. In the usual metric sense an arclength-parametrized Teichm\"uller geodesic satisfies $d_{\mathcal T}(\gamma(s),\gamma(t))=|s-t|$. The record does not specify the parameter interval, whether a complete line or a ray is intended, or whether only an unparametrized image is requested.

For matrix or operator period coordinates, ``imaginary axis'' could mean different subsets, and the source does not provide a one-parameter expression or explain how such an axis lies in the locus. The question asks for an image of a geodesic, not that the period map be an isometry or that every path in a purely imaginary subspace lift geodesically. A precise statement requires recovery of the special construction and its coordinates.
\subsection{Short English statement}
The question asks whether a particular imaginary axis in an infinite-genus period-matrix construction is exactly the period image of a Teichm\"uller geodesic. The locus, axis, surface, period map, and geodesic range are undefined in the supplied fragment.
\subsection{Sources}
\begin{itemize}
\item Local UnsolvedMath registry, permanent ID 30000508, code catalogue entry, ``Imaginary-Axis Geodesics in Period-Matrix Loci.'' Proposed year 2006; proposer unspecified.
\item \emph{Teichm\"uller Space (Classical and Quantum)} (2007), Oberwolfach Reports, pp.\ 1537--1614, DOI 10.4171/owr/2006/26; \url{https://doi.org/10.4171/owr/2006/26}. The displayed report year 2007 differs from the DOI and proposed year 2006.
\item Statement status corrected\_verified; the supplied verification says official-report content was checked and notation, context, and question form normalized without changing the source claim.
\item Literature triage checked 2026-08-21: open. It describes a special infinite-genus construction, says the fragment is not a complete theorem, and reports no located source settling its general assertion. Created 2026-08-25; updated 2026-09-16. Time-sensitive status and source claims require independent verification.
\item Registry curation by Ulam AI / UnsolvedMath Contributors, CC BY 4.0; the underlying report retains its supplied provenance. Unreviewed statement-only draft.
\end{itemize}
\end{document}
